\subsection{The pointed Heisenberg covers and the exceptional prime two}

\begin{result}{C.1}{Proved}{Geometry of every odd-prime Heisenberg cover}
Let \(\ell\) be odd and \(t_0\in U(\mathbb Q)\).
Let \(H_\ell\) be the universal two-generator quotient of class at
most two and exponent dividing \(\ell\). It is the group of upper
unitriangular \(3\times3\) matrices over \(\mathbb F_\ell\), of
order \(\ell^3\). The inverse image of the point section at \(t_0\)
in the arithmetic \(H_\ell\)-quotient defines a connected finite
etale cover \(Y_{\ell,t_0}\to U\) of degree \(\ell^3\), with a
specified rational point above \(t_0\). Its compactification has
\[
 g(\overline Y_{\ell,t_0})=1+\frac{\ell^2(\ell-3)}2,
 \qquad |\overline Y_{\ell,t_0}\setminus Y_{\ell,t_0}|=3\ell^2.
\]
Its central quotient is the pointed diagonal curve
\[
 C_{\ell,t_0}:\quad t_0 u^\ell+(1-t_0)v^\ell=1,
 \qquad (u,v)=(1,1),
\]
and \(\overline Y_{\ell,t_0}\to\overline C_{\ell,t_0}\)
is an unramified cover of degree \(\ell\).
Thus the genus-one case \(\ell=3\) is exceptional; the genera at
\(\ell=5,7\) are already \(26,99\).
\end{result}
\begin{proof}
Write matrices as \((a,b,c)\), with multiplication
\[
 (a,b,c)(a',b',c')=(a+a',b+b',c+c'+ab').
\]
The commutator of \(x=(1,0,0)\) and \(y=(0,1,0)\) generates
the central subgroup \(\{(0,0,c)\}\). Collection gives at most
\(\ell^3\) elements in any two-generator class-two exponent-\(\ell\)
group, and these matrices attain that bound. Since
\((a,b,c)^n=(na,nb,nc+\binom n2ab)\), their exponent divides
\(\ell\) for odd \(\ell\). The quotient is characteristic in the
geometric group, so the arithmetic action descends. Taking the inverse
image of the specified arithmetic section is the usual Galois-category
construction of the stated pointed cover. Its geometric subgroup is the
kernel of the epimorphism onto \(H_\ell\); therefore it is connected
and of the stated degree.

The three geometric branch inertia generators are \(x,y,(xy)^{-1}\).
Each has order \(\ell\), and each has nonzero image in
\(H_\ell^{\rm ab}\). Above each of \(0,1,\infty\) there are
\(\ell^2\) points, all of ramification index \(\ell\).
Riemann--Hurwitz gives
\[
 2g-2=-2\ell^3+3\ell^3(1-1/\ell)=\ell^2(\ell-3).
\]
The quotient by the center has the two Kummer equations
\(t/t_0=u^\ell\) and \((1-t)/(1-t_0)=v^\ell\).
Their elimination gives the displayed smooth diagonal curve.
The center meets each branch inertia group trivially, because an
element of a cyclic inertia group with trivial abelian image must be
its identity. Hence the intermediate central cover is unramified even
over the boundary. Its Riemann--Hurwitz formula agrees with the
genus \((\ell-1)(\ell-2)/2\) of the diagonal curve.
These arguments take place over \(\overline{\mathbb Q}\); pointing
supplies the asserted descent and rational point over \(\mathbb Q\).
\end{proof}

\begin{result}{C.2}{Refuted}{The exponent-prime degree formula cannot include two unchanged}
There is no nonabelian group of exponent two. In particular the
universal rank-two class-two exponent-two quotient has order four,
not eight. If ``Heisenberg at two'' instead means
\(\mathrm U_3(\mathbb F_2)\), that group has exponent four and is
dihedral of order eight. A cover with branch generators its two standard
involutions has branch orders \((2,2,4)\), genus zero, and ten boundary
points. These are two different conventions, and neither is obtained
by substituting two into the odd-prime exponent-\(\ell\) assertion.
The exponent-two quotient instead has genus zero and six boundary points.
A pointed degree-eight dihedral model above \(t_0=3\) is explicit:
\[
\begin{gathered}
 a(s)=\frac{1+6s^2}{1-6s^2},\qquad
 b(s)=\frac{2s}{1-6s^2},\\
 u(s)=a(s)^2+6b(s)^2+4a(s)b(s),\qquad
 v(s)=a(s)^2+6b(s)^2+6a(s)b(s),\\
 t(s)=3u(s)^2,\qquad 1-t(s)=-2v(s)^2.
\end{gathered}
\]
The point \(s=0\) lies over \(t=3\). The open source is
\(\mathbb P^1_s\) with the ten points above \(0,1,\infty\) removed.
\end{result}
\begin{proof}
In a group of exponent two,
\(xy=(xy)^{-1}=y^{-1}x^{-1}=yx\). The universal two-generator
such group is therefore \((\mathbb Z/2)^2\).
In the matrix group above at \(\ell=2\), \(x^2=y^2=1\), but
\((xy)^2=(0,0,1)\ne1\), so \(xy\) has order four. These relations
give the dihedral group. A geometric regular cover with this specified
monodromy has \(4,4,2\) boundary points. Riemann--Hurwitz is
\(2g-2=-16+8(1/2+1/2+3/4)=-2\).
This calculation does not silently identify this labelled quotient with
the characteristic exponent-two quotient. The finite group calculations
are independently enumerated by \texttt{code/root\_heisenberg.py}.
For the exponent-two quotient the three inertia orders are all two,
and the same formula at degree four gives genus zero and six punctures.

For the pointed model, \(a^2-6b^2=1\), and direct multiplication in
\(\mathbb Q(\sqrt6)\) gives
\[
 3u+v\sqrt6=(3+\sqrt6)(a+b\sqrt6)^2.
\]
Taking norms gives \(3u^2-2v^2=1\). Over the algebraic closure put
\(w=a+b\sqrt6\). The intermediate cover is the square map
\(w\mapsto w^2\) on the norm-one torus, followed by the degree-four
Kummer map of \(3u^2-2v^2=1\). The two sign changes of \(u,v\)
act on \(z=w^2\) as \(z\mapsto c/z\) and
\(z\mapsto-c/z\), for a nonzero constant \(c\).
Their lifts are involutions \(w\mapsto\alpha/w\) and
\(w\mapsto\beta/w\) with \(\alpha^2=-\beta^2\).
Their product has order four. They generate a geometric dihedral
group of order eight, equal to the degree of the displayed map.
The four simple zeros of the numerator of \(u\), the four of
\(v\), and the two poles at \(1-6s^2=0\) give branch orders
\((2,2,4)\). Nonvanishing discriminants and pairwise resultants
are checked exactly in \texttt{code/root\_cover\_checks.py}.
Their total ramification is \(14=2\cdot8-2\), so there is no
additional branching. Since \(u(0)=v(0)=1\), this is the
pointed form at three. The geometric quotient with labelled inertia
generators two independent involutions in the dihedral group is unique
up to an automorphism of that group. It is stable under the actual
tangential arithmetic action: the two meridians have odd cyclotomic
powers, and the intervening path commutator becomes central at class
two. To identify the arithmetic twist, let \(N\) be this stable
geometric kernel and \(H\) the arithmetic subgroup of the displayed
cover. Then \(H/N\) maps isomorphically to \(G_{\mathbb Q}\).
The rational point \(s=0\) above three puts the point-three section
inside \(H/N\); two sections with one contained in the other are equal.
Thus this is exactly the point-three twist of the indicated arithmetic
quotient, with exponent four explicitly retained.
\end{proof}

\begin{result}{C.3}{Proved}{Passing to a cover preserves the appropriate prime, not every prime}
Let \(Y\to U\) be a connected geometrically Galois cover whose
geometric group is a nontrivial \(p\)-group of order \(d\).
If a geometrically pro-\(p\) section factors through its associated
open arithmetic subgroup, that subgroup is the full geometrically
pro-\(p\) arithmetic fundamental group of \(Y\).
For the full pronilpotent group of \(U\), the induced geometric open
subgroup is instead
\[
 K_p\times\prod_{\ell\ne p}F_\ell,
 \qquad [F_p:K_p]=d.
\]
It is not the full pronilpotent geometric fundamental group of \(Y\).
For \(\ell\ne p\), the latter has free pro-\(\ell\) rank \(d+1\),
whereas the inherited factor has rank two.
\end{result}
\begin{proof}
Put \(K=\pi_1(Y_{\overline{\mathbb Q}})\subset\Pi\).
If \(V\triangleleft K\) has \(p\)-power index, its intersection
over the finitely many cosets of \(K\) in \(\Pi\) is normal in
\(\Pi\) and has \(p\)-power index. This intersection is contained
in \(V\), proving that the topology induced from the pro-\(p\)
completion of \(\Pi\) is the full pro-\(p\) topology on \(K\).
The completion therefore identifies \(K^{(p)}\) with \(K_p\).
The arithmetic assertion follows by the same kernel comparison over
\(G_{\mathbb Q}\), with the specified section retained.

In the pronilpotent product, the finite \(p\)-group quotient kills
every factor of characteristic different from \(p\). Its kernel is
exactly the displayed product. On the other hand, the Euler
characteristic of \(U\) is \(-1\), so the geometric free group of
the degree-\(d\) affine cover has rank \(1+d\).
Its pro-\(\ell\) abelianization consequently has rank \(d+1\) over
\(\mathbb Z_\ell\). The inherited factor has abelianization
\(\mathbb Z_\ell^2\); these ranks disagree. Thus the natural
comparison map from the cover's full pronilpotent group to the inherited
subgroup cannot be an isomorphism. In particular, a section of the
inherited subgroup is not already a section of the former group.
This is the precise missing comparison in a proposed application of
a full section theorem on a cover to a downstairs pronilpotent section.
\end{proof}

\begin{result}{C.4}{Proved}{Compact finite separation and the limit of the pentagon analogy}
Let \(X\) be a compact space, \(H\) a profinite group, and
\(f:X\to H\) continuous. If \(1\notin f(X)\), there is an open
normal subgroup \(V\triangleleft H\) with \(f(X)\cap V=\varnothing\).
Thus an exclusion proved on an actual compact family is detected in
one finite quotient. This general fact proves the topological inference
used in GT\_TOWER\_CHECK A.4.1 conditional on that work's arithmetic
vanishing input; it does not verify A.7 or the explicit compiler A.4.57.
\end{result}
\begin{proof}
The set \(f(X)\) is compact and hence closed in the Hausdorff group
\(H\). Its complement is an open neighbourhood of one. Open normal
subgroups form a neighbourhood basis in a profinite group, so one of
them lies in this complement. The assertion about the finite quotient
is exactly the assertion that the identity coset is avoided in \(H/V\).
For the varying ghost families, A.101 disproves the required compactness.
Consequently this lemma cannot be applied to those families without
an additional compact parameter space that actually contains them.
\end{proof}

\begin{result}{C.5}{Known}{Why the modular Fermat theorem does not descend to this quotient}
\citet[Theorem A and Theorem 5.14]{Steklov2026} proves a cuspidality
theorem for full Selmer sections of affine Fermat curves of prime
exponent at least seven. The residual Legendre representations used
there do not factor geometrically through the pronilpotent quotient of
the tripod at an odd coefficient prime \(\ell\ge5\).
Thus this theorem is not a proof of the present main assertion.
\end{result}
\begin{proof}
The geometric Legendre monodromy is generated, up to the choice of
sign of one meridian, by
\[
 \begin{pmatrix}1&2\\0&1\end{pmatrix},\qquad
 \begin{pmatrix}1&0\\-2&1\end{pmatrix}.
\]
The precise monodromy calculation is \citet[\S5.1]{Stix2015}.
Modulo an odd \(\ell\), their powers give every upper and lower
elementary matrix, because two is invertible. Elementary row reduction
then gives all of \(\mathrm{SL}_2(\mathbb F_\ell)\).
Its upper unipotent Sylow \(\ell\)-subgroup is conjugated to the
distinct lower one, so it is not normal. A finite nilpotent group has
normal Sylow subgroups; this image is therefore not nilpotent.
Every finite image of a pronilpotent group is nilpotent, which proves
the asserted nonfactorization. There is a separate issue as well:
ordinary abelian local data do not establish full Selmer local
geometricity on the Fermat cover. The cover comparison in C.3 shows
exactly why passage to that cover cannot supply the missing quotient.
\end{proof}

\begin{result}{C.6}{Proved}{Fixed-coordinate reciprocity does not exclude the constructed ghosts}
Let \(A\in\widehat{\mathbb Q^\times}\), and let
\(\theta:G_{\mathbb Q}\to\mathbb Q/\mathbb Z\) be a finite
character of order dividing \(n\). The local reciprocity sum for the
local Kummer coordinates of \(A\), paired with \(\theta\), is zero.
It is well defined even when the ordinary local coordinates of \(A\)
do not form an idele. The same holds for \(B\) and for every fixed
monomial in \(A,B\). Hence reciprocity on these base unit coordinates
alone cannot distinguish a ghost from a global completed pair.
\end{result}
\begin{proof}
Choose a rational representative \(a_n\in\mathbb Q^\times\) of
the image of \(A\) in \(\mathbb Q^\times/(\mathbb Q^\times)^n\).
At every place the local pairing depends only on this image, so it is
the same pairing as for \(a_n\). Only finitely many local terms for
the rational element are nonzero. Global reciprocity says their sum,
including the real term, is zero
\citep[VII, Theorem 8.1, p.~221]{RootMilneCFT}. Independence of the representative
follows because an \(n\)-th power pairs trivially with \(\theta\).
The argument applies verbatim to a fixed monomial. It uses finite
characters throughout and does not assert that an infinite-support
family is an element of the restricted idele group.
\end{proof}

\begin{result}{C.7}{Open}{The remaining uniform-target and countersign assertions}
No object excluding every ordinary infinite-support pair has been
identified here, and no theorem here excludes all the proposed types of
object. The all-prime genera in C.1 do not compute the Selmer groups of
their Jacobians, or supply full sections on them. GT\_TOWER\_CHECK
C.3.5 asserts passage of the displayed compact point-order tests only;
it does not assert passage of every pro-three quotient. Existence of a
full section for every remaining member of that family is undecided.
The stronger countersign condition involving all relevant covers is
also undecided; C.6 refutes only the interpretation using fixed base
unit coordinates. These are the uncompleted parts of C1--C4 and
Leads L3, L5, L6.
\end{result}

\begin{result}{C.8}{Proved}{The missing pointwise hypothesis in an all-prime pointed-cover test}
Fix \(t_0\in U(\mathbb Q)\). Suppose an ordinary completed pair
\((A,B)\) has, for every prime \(\ell\), the same reduction modulo
\(\ell\)-th powers as \((t_0,1-t_0)\). Then
\((A,B)=(\kappa(t_0),\kappa(1-t_0))\).
In particular, factorization of a single section through the
point-\(t_0\) Heisenberg cover at every prime would already force its
ordinary abelian pair to be rational. Such factorization is not an
automatic property of a full section: the rational-point section at
two does not factor through the point-three Heisenberg cover at five.
\end{result}
\begin{proof}
At any fixed prime \(q\), the difference between \(v_q(A)\) and
\(v_q(t_0)\) is an integer, by ordinarity. The hypothesis makes it
divisible by every prime \(\ell\), hence zero. The same argument
applies to \(B\). The sign coordinate is fixed by the reduction
modulo two. The description of the power completion in E.101 then
gives equality of both global completed classes, without passing to
an unproved closure of rational points.

A lift to the specified pointed Heisenberg cover has the same section
class in its finite geometric quotient. Composing with the
abelianization of that quotient gives the hypothesis just used.
For the counterexample, the first coordinate of the point-two section
has valuation one at two; the first coordinate of point three has
valuation zero there. These are different modulo five. The sections
therefore cannot be conjugate even in the abelianization of the
degree-\(125\) Heisenberg quotient. Geometric conjugation does not
change their abelian cohomology classes. Thus the former section
cannot factor through the latter pointed cover. This gives the exact
missing implication in an attempted use of all point-three covers as
universal section neighbourhoods.
\end{proof}
